\documentclass[11pt,a4paper]{article}
% Common preamble for the hint sheets (\input from each file)
\usepackage[utf8]{inputenc}
\usepackage[T1]{fontenc}
\usepackage{lmodern}
\usepackage{amsmath,amssymb,amsthm,mathtools}
\usepackage[a4paper,margin=2.3cm]{geometry}
\usepackage{enumitem}
\usepackage{xcolor}
\usepackage[hidelinks]{hyperref}
\renewcommand{\leq}{\leqslant}
\renewcommand{\le}{\leqslant}
\renewcommand{\geq}{\geqslant}
\renewcommand{\ge}{\geqslant}
\setlength{\parindent}{0pt}
\setlength{\parskip}{0.35em}
% \hintstitle{Sheet title}{Class work / Homework}
\newcommand{\hintstitle}[2]{\begin{center}{\LARGE\bfseries Hints: #1\par}\vspace{0.5em}{\large #2\par}\end{center}
\noindent\textit{The hints to each problem are cumulative. Read one, try again, and open the next only if you are still stuck.}\par\medskip\hrule\medskip}
\newcommand{\problem}[1]{\par\bigskip{\large\bfseries Problem #1}\par\nopagebreak\smallskip}
\newcommand{\subpart}[1]{\par\smallskip\textbf{(#1)}\par\nopagebreak}
\newcommand{\hint}[1]{\par\smallskip\noindent\textcolor{black!70}{\textbf{Hint #1.}}\enspace}
\newcommand{\note}[1]{\par\smallskip\noindent\textit{Note.} #1\par}
\providecommand{\lt}{<}
\providecommand{\gt}{>}
\providecommand{\tightlist}{\setlength{\itemsep}{0pt}\setlength{\parskip}{0pt}}
\emergencystretch=3em

\begin{document}
\hintstitle{Methods of proof -- 2}{Homework}

\problem{1}
\hint{1} Lines: the $n$-th line meets the previous $n-1$ lines in $n-1$ distinct points. These points cut it into $n$ pieces, and each piece splits one old region into two. So $f(n)=f(n-1)+n$ with $f(0)=1$, which gives $f(n)=1+\frac{n(n+1)}2=\frac{n^2+n+2}2$.
\hint{2} In $\mathbb R^d$, let $f(n,d)$ be the number of regions cut out by $n$ hyperplanes in general position. The $n$-th hyperplane $H\cong\mathbb R^{d-1}$ is cut by the others into $f(n-1,d-1)$ regions, because their traces on $H$ are $n-1$ hyperplanes of $H$ in general position. Each such region splits one old region.
\hint{3} So $f(n,d)=f(n-1,d)+f(n-1,d-1)$, with $f(0,d)=1$ and $f(n,0)=1$. This is Pascal's rule, and the double induction on $(n,d)$ (e.g.\ on $n+d$) gives $f(n,d)=\sum_{i=0}^d\binom ni$.

\problem{2}
\hint{1} Let $\Delta P(x)=P(x+1)-P(x)$. If $\deg P=n$ and $P$ is integer at $a,a+1,\dots,a+n$, then $\Delta P$ has degree $n-1$ and is integer at $a,\dots,a+n-1$. Induct on the degree: all $\Delta^kP(a)$, $k=0,\dots,n$, are integers, and $\Delta^{n+1}P=0$.
\hint{2} Newton's forward-difference formula: $P(x)=\sum_{k=0}^n\Delta^kP(a)\binom{x-a}k$. Prove it by induction on $n$ too, applying it to $\Delta P$ and using $\Delta\binom{x-a}{k}=\binom{x-a}{k-1}$. Hence $P$ is integer at every integer.
\hint{3} The same argument with $a=0$ shows that every integer-valued polynomial is an integer combination of the $\binom xk$. The $\binom xk$ are integer-valued and have distinct degrees, so they are linearly independent. Therefore they form a basis of the $\mathbb Z$-module of integer-valued polynomials.

\problem{3}
\hint{1} Cauchy--Schwarz: the case of two terms, $(ab+cd)^2\le(a^2+c^2)(b^2+d^2)$, is equivalent to $(ad-bc)^2\ge0$. For the step, put $A=\big(\sum_{i\le n}a_i^2\big)^{1/2}$ and $B=\big(\sum_{i\le n}b_i^2\big)^{1/2}$. Then $\sum_{i\le n+1}a_ib_i\le AB+a_{n+1}b_{n+1}$ by the induction hypothesis. Apply the two-term case to $(A,a_{n+1})$ and $(B,b_{n+1})$.
\hint{2} AM--GM forward: the case $n=2$ is $(\sqrt x-\sqrt y)^2\ge0$. From $n$ to $2n$, apply the case $n$ to each half and then the case $2$ to the two means. This gives all powers of $2$.
\hint{3} Backward step $n\to n-1$: given $x_1,\dots,x_{n-1}$ with mean $m$, apply the $n$-term inequality to $x_1,\dots,x_{n-1},m$. This gives $m^n\ge x_1\cdots x_{n-1}\,m$, hence $m^{n-1}\ge x_1\cdots x_{n-1}$.
\hint{4} Why the scheme: a direct step $n\to n+1$ is awkward, because the mean of $n+1$ numbers does not split into the mean of $n$ of them plus one term in a way that AM--GM for $n$ handles directly. Doubling and adding the mean as an extra variable are the natural moves. Together they reach every $n$, since every $n$ lies below some power of $2$.

\problem{4}
\hint{1} Class equation: $|G|=|Z(G)|+\sum_i[G:C_G(x_i)]$, where the sum runs over representatives of the conjugacy classes with more than one element. Each such index divides $p^k$ and is greater than $1$, so it is divisible by $p$.
\hint{2} Hence $p\mid|Z(G)|$, and since $e\in Z(G)$, the centre is non-trivial. An induction on $k$ is not really needed here. If you want one, use it for the corollary that a $p$-group has normal subgroups of every order: pass to $G/Z$, of order $p^{k'}$ with $k'<k$.
\hint{3} Order $p^2$: if $Z\neq G$, then $|Z|=p$ and $G/Z$ has order $p$, so it is cyclic. If $G/Z$ is cyclic, generated by $gZ$, then every element is $g^iz$ with $z\in Z$, and any two such elements commute. So $G$ is abelian, which contradicts $Z\neq G$.

\problem{5}
\subpart{$x^2+y^2+z^2=2xyz$}
\hint{1} Strengthen: show that $x^2+y^2+z^2=2^m xyz$ has no positive solutions for any $m\ge1$. Look at the equation modulo $4$, recalling that squares are $0$ or $1\pmod 4$.
\hint{2} The right side is even, so an even number of $x,y,z$ are odd. If exactly two are odd, the left side is $\equiv2\pmod4$, while the right side is divisible by $4$ (it has a factor $2$ and one even variable). So all three are even.
\hint{3} Substituting $x=2x'$, $y=2y'$, $z=2z'$ gives $x'^2+y'^2+z'^2=2^{m+1}x'y'z'$, a solution of the same type with a smaller $x+y+z$. This is an infinite descent, so no solution exists.
\subpart{$x^3+2y^3=4z^3$}
\hint{1} Assume a non-trivial integer solution with $|x|+|y|+|z|$ minimal (the hint covers integers, not just positive ones). Since $x^3$ is even, $x$ is even.
\hint{2} With $x=2x_1$: $8x_1^3+2y^3=4z^3$, so $4x_1^3+y^3=2z^3$, and $y$ is even. With $y=2y_1$: $2x_1^3+4y_1^3=z^3$, and $z$ is even. With $z=2z_1$: $x_1^3+2y_1^3=4z_1^3$, a smaller solution.

\problem{6}
\hint{1} Equivalence $A\sim PAQ$ with $P,Q$ invertible means doing row and column operations. If $A\neq0$, move a non-zero entry to position $(1,1)$, scale it to $1$, and clear the rest of the first row and column.
\hint{2} This gives $\begin{pmatrix}1&0\\0&A'\end{pmatrix}$ with $\operatorname{rank}A'=r-1$. Apply the induction hypothesis on the rank to $A'$.
\hint{3} Nilpotent $N$ ($N^k=0$): $N$ is not invertible, so there is a $v\neq0$ with $Nv=0$. In a basis starting with $v$, $N=\begin{pmatrix}0&*\\0&N'\end{pmatrix}$, and $N'$ is nilpotent because the block-triangular form gives $(N')^k$ as the lower-right block of $N^k=0$. Make $N'$ strictly upper triangular by induction on the size. Conjugating by the matrix $\operatorname{diag}(1,S)$ keeps the first column zero.

\problem{7}
\note{``Cannot be proved by induction on any natural parameter'' is not a precise mathematical statement. The hints read it as: no argument that works only with finitely many points or steps can prove it. Then they give the compactness proof and the rule.}
\hint{1} Compare with sets where the statement fails: $f(x)=1/x$ on $(0,1]$, and $f(x)=1/(x-1/\sqrt2)$ on $[0,1]\cap\mathbb Q$. Both are continuous, and every finite configuration of points behaves exactly as on $[0,1]$. So any proof must use completeness of $[0,1]$, which an induction over finitely many steps never sees.
\hint{2} Compactness proof: if $f$ is unbounded, choose $x_n$ with $|f(x_n)|>n$. By Bolzano--Weierstrass, a subsequence converges to some $x\in[0,1]$, and continuity at $x$ gives a contradiction. Alternatively, let $s=\sup\{t:f$ is bounded on $[0,t]\}$. Show $s>0$, that $f$ is bounded on $[0,s]$, and that $s=1$ (``real induction'').
\hint{3} The rule: induction on $\mathbb N$ is available when the objects carry a natural-number parameter, and each object can be reduced to finitely many objects with smaller parameter. Behind this is well-foundedness. Over the reals, where descending sequences need not stop, it is replaced by sup, compactness, or a limit argument.

\problem{8}
\hint{1} First prove Euler's formula $V-E+F=2$ for connected plane graphs by induction on the number of edges. Removing an edge on a cycle merges two faces. If there is no cycle, the graph is a tree. Count edge--face incidences to get $E\le3V-6$ for $V\ge3$.
\hint{2} The sum of degrees is $2E\le6V-12<6V$, so some vertex has degree at most $5$. This property passes to subgraphs, which is what the induction on the number of vertices uses.
\hint{3} Five colours: remove a vertex $v$ of degree $\le5$ and colour the rest by induction. If the neighbours use at most $4$ colours, colour $v$. Otherwise $v$ has five neighbours $v_1,\dots,v_5$ in cyclic order, with five different colours.
\hint{4} Kempe chains: consider the subgraph spanned by colours $1$ and $3$. If $v_1$ and $v_3$ lie in different components of it, swap the two colours in the component of $v_1$, which frees colour $1$ for $v$. Otherwise a $1$--$3$ path joins $v_1$ to $v_3$. Together with $v$ it forms a closed curve that separates $v_2$ from $v_4$, so the $2$--$4$ swap works instead.

\problem{9}
\hint{1} Maximal ideal: apply Zorn's lemma to the proper ideals, ordered by inclusion (the zero ideal shows the family is non-empty). The union of a chain of ideals is an ideal, and it is proper because it does not contain $1$. This is exactly where the unit is used.
\hint{2} Maximal independent set: the same argument, as in the class work. The union of a chain of independent sets is independent, since linear dependence involves only finitely many vectors.
\hint{3} Countably generated $V=\operatorname{span}\{v_1,v_2,\dots\}$: go through $v_1,v_2,\dots$ and keep $v_n$ if it is not in the span of the vectors kept so far. By induction the kept vectors are independent and span $\operatorname{span}\{v_1,\dots,v_n\}$ at stage $n$. So they form a basis, and no choice principle beyond ordinary induction is needed.

\problem{10}
\hint{1} Sperner in dimension $n$: take a triangulation of the simplex $[e_0,\dots,e_n]$. Each vertex gets a label $i$ from the coordinates that are positive at that vertex (i.e.\ from the indices of the smallest face containing it). Claim: the number of fully labelled small simplices is odd. Dimension $1$: along the segment, the label changes from $0$ to $1$ an odd number of times.
\hint{2} Induction step (door counting): call an $(n-1)$-face labelled $0,\dots,n-1$ a door. Count pairs (small simplex, door on its boundary). A fully labelled simplex has exactly one door, any other simplex has $0$ or $2$, and an interior door is counted twice.
\hint{3} The doors on the boundary all lie on the face $[e_0,\dots,e_{n-1}]$, and by the induction hypothesis there are an odd number of them. Hence the number of fully labelled simplices is odd, in particular non-zero.
\hint{4} Brouwer: for continuous $f:\Delta\to\Delta$, label a vertex $v$ by some $i$ with $v_i>0$ and $f(v)_i\le v_i$. Such an $i$ exists, since the coordinates of both $v$ and $f(v)$ sum to $1$. The labelling is a Sperner labelling. Take finer and finer triangulations, and use compactness to get a limit point $x$ of fully labelled simplices. Then $f(x)_i\le x_i$ for all $i$, so $f(x)=x$. For $n=1$ this is the intermediate value theorem applied to $f(x)-x$.

\end{document}
